\section{The definite integral}
\begin{itemize}

    \item Denoted by 
	    \begin{equation*}
	        \int^{b}_{a} f(x) \; dx
	    \end{equation*}
    \item  
     Geometric definition of Definite integral 
     \begin{figure}[h]
	\centering
	\def\svgwidth{0.7\columnwidth}
	\import{./images/}{definite_integral.pdf_tex}
     \end{figure}
     Area of the region above the $x$ axis and below the curve $y = f(x)$
     between two vertical lines $x=a$ and $x=b$ . $a$ and $b$ are 
     upper and lower bound respectively. 
\end{itemize}
Summation Notation
\begin{equation*}
    \sum_{i=1}^{a} = a_{1} + a_{2} + .... + a_{n}
\end{equation*}

\subsection{Riemann Sums}

steps to evaluate $\int^{b}_{a} f(x) \; dx$ using Riemann sums 

\begin{enumerate}

    \item Devide $[a,b]$ into $n$ equal sub intervals 
    \begin{figure}[H]
        \centering
	\def\svgwidth{0.5\columnwidth}
	\import{./images/}{riemann1.pdf_tex}
    \end{figure}
    \item choose a point $c_{i}$ with in the $i^{th}$ subinterval 
	 and choose $f(c_{i})$ be the height of the $i^{th}$ rectangle. 
    \begin{figure}[H]
        \centering
	\def\svgwidth{0.5\columnwidth}
	\import{./images/}{riemann2.pdf_tex}
    \end{figure}
    \item Add up the areas of the $n$ rectangles. the  total areas 
      of n rectangles is 
      \begin{equation*}
	      f(c_{1}) \Delta x + f(c_{2}) \Delta x + .....
	      + f(c_{n}) \Delta x = \sum_{i=1}^{n}f(c_{i}) \Delta x
      \end{equation*}
      \item Take the limit as the rectangles becomes infintesimally thin
	      ($\Delta x \to 0$ or equivalent $n \to \infty$) . The limit
	is actual area under curve between $a$ and $b$ 
	\begin{equation*}
	    \lim_{n \to \infty} \sum_{i= 1}^{n} f(c_{i}) \Delta x
	    = \int^{b}_{a} f(x) dx
	\end{equation*}
\end{enumerate}
note : if we pick $c_{i}$ to be left end point then the Riemann sum 
is left Riemann sum , if $c_{i}$ is right end point then Riemann sum is Tight Riemann sum.

However in the limit $n\to \infty$ (so that $\Delta x \to 0$) , The 
distinction between these two sums is not longer needed , The limit 
of any Riemann sum , no matter what $c_{i}$ with in subintervals are is 
equal to exact area under curve. 
    
\subsection{Cummulative Sums}
For instance if $x$ represent time (hr) and f(x) denotes velocity 
(km/hr). thus the $\int^{b}_{a} dx$ represent distance (km).
